**Riesz--Thorin interpolation theorem.** Let \((X,\mu)\) and \((Y,\nu)\) be \(\sigma\)-finite measure spaces, and let \(1 \leq p_0,p_1,q_0,q_1 \leq \infty\). For \(0<\theta<1\), define \(p_\theta\) and \(q_\theta\) by
\[
\frac{1}{p_\theta}
=
\frac{1-\theta}{p_0}
+
\frac{\theta}{p_1},
\qquad
\frac{1}{q_\theta}
=
\frac{1-\theta}{q_0}
+
\frac{\theta}{q_1},
\]
with the convention that \(1/\infty=0\). If \(T\) is a linear operator such that
\[
T:L^{p_0}(X)\to L^{q_0}(Y),
\qquad
T:L^{p_1}(X)\to L^{q_1}(Y)
\]
are bounded with
\[
\|T\|_{L^{p_0}\to L^{q_0}} \leq M_0,
\qquad
\|T\|_{L^{p_1}\to L^{q_1}} \leq M_1,
\]
then
\[
T:L^{p_\theta}(X)\to L^{q_\theta}(Y)
\]
is bounded and
\[
\|T\|_{L^{p_\theta}\to L^{q_\theta}}
\leq
M_0^{1-\theta}M_1^\theta .
\]
This is the standard \(L^p\)-operator norm form of the Riesz--Thorin interpolation theorem. ([en.wikipedia.org](https://en.wikipedia.org/wiki/Riesz%E2%80%93Thorin_theorem))

Now regard the matrix \(H\) as a linear operator on \(\ell_p^m=L^p(\{1,\dots,m\})\) with counting measure. We are given that
\[
\|H\|_{2\to 2}=\sqrt{m},
\qquad
\|H\|_{\infty\to\infty}=m.
\]

First consider \(r=2\). Then
\[
\|H\|_{2\to 2}
=
\sqrt{m}
=
m^{1/2}
=
m^{1-1/2},
\]
so the desired bound holds for \(r=2\).

Now let \(r>2\). Apply Riesz--Thorin with
\[
p_0=q_0=2,
\qquad
p_1=q_1=\infty,
\qquad
M_0=\sqrt{m},
\qquad
M_1=m.
\]
Choose \(\theta\in(0,1)\) by
\[
\frac{1}{r}
=
\frac{1-\theta}{2}
+
\frac{\theta}{\infty}
=
\frac{1-\theta}{2}.
\]
Since \(1/\infty=0\), this gives
\[
1-\theta=\frac{2}{r},
\qquad
\theta=1-\frac{2}{r}.
\]
Because \(r>2\), indeed \(\theta\in(0,1)\). With this choice,
\[
p_\theta=q_\theta=r.
\]
Therefore Riesz--Thorin gives
\[
\|H\|_{r\to r}
\leq
\|H\|_{2\to 2}^{1-\theta}
\|H\|_{\infty\to\infty}^{\theta}
=
(\sqrt{m})^{1-\theta}m^\theta.
\]
Using \(1-\theta=2/r\) and \(\theta=1-2/r\), we obtain
\[
(\sqrt{m})^{1-\theta}m^\theta
=
(\sqrt{m})^{2/r}m^{1-2/r}
=
m^{1/r}m^{1-2/r}
=
m^{1-1/r}.
\]
Hence
\[
\boxed{\|H\|_{r\to r}\leq m^{1-1/r}}
\]
for every real \(r\geq 2\).